% Section 6: Synthesis (~1,200 words)

\subsection{Four Properties of a Valid Redistricting Algorithm}

Drawing together the five empirical papers, this section articulates
four properties that the C-track collectively validates for METIS
recursive bisection redistricting.

\subsubsection{Property 1: Spatial Robustness}

\begin{quote}
\textit{The algorithm produces consistent outputs regardless of the
spatial resolution of the input units, within a 130$\times$ range
of unit counts.}
\end{quote}

Evidence: C.1 demonstrates that Polsby-Popper compactness varies by
$\Delta PP < 0.005$ (mean) and $< 0.011$ (maximum, Nevada) across
tract, block-group, and block resolutions. VRA-relevant majority-minority
district counts vary by $\pm 1$ nationally across the same resolution
range.

Interpretation: The algorithm responds to the underlying geographic
structure (terrain, settlement patterns, administrative boundaries)
rather than to the choice of building block. This is a prerequisite
for algorithmic legitimacy: a method whose conclusions depend on an
arbitrary implementation choice (tract vs.\ block) would not be
trustworthy.

\subsubsection{Property 2: Temporal Stability}

\begin{quote}
\textit{The algorithm produces stable outputs across the three
available census cycles (2000, 2010, 2020), with geographic factors
explaining $3.2\times$ more variance than temporal factors.}
\end{quote}

Evidence: C.2's variance decomposition finds $\sigma^2_\text{geo} /
\sigma^2_\text{temp} = 3.2$. Within-slice cross-census correlations
have median $r = 0.73$. National mean PP varies by approximately 10\%
over twenty years (0.412 to 0.441).

C.3 demonstrates that recursive bisection provides $0.8$ pp better
decadal boundary stability than n-way partitioning---a structural
advantage from hierarchical representation of geographic scale.

Interpretation: An algorithm validated in 2020 does not require
complete revalidation for 2030. The geographic basis of its performance
is stable; only the input data changes. This property supports the DIA's
ten-year certification framework.

\subsubsection{Property 3: Demographic Neutrality}

\begin{quote}
\textit{The algorithm's performance (compactness, VRA compliance)
is driven by geographic structure, not demographic composition.
Demographic shifts between census years do not systematically bias
the algorithm toward or away from any population group.}
\end{quote}

Evidence: C.2 demonstrates that temporal variance ($\sigma^2_\text{temp}$)
is small relative to geographic variance ($\sigma^2_\text{geo}$). Slices
with high demographic change (Nevada, Arizona, Florida) show higher
temporal variance, but even these remain within the distribution of
cross-geographic variance.

C.4 shows that the geographic clustering of states by compactness
performance is stable: 47 of 50 states remain in the same cluster
across 2000--2020, despite substantial demographic change over that
period.

Interpretation: The algorithm does not develop differential performance
for states with changing racial compositions, political demographics,
or urbanization levels. Its behavior is \emph{geography-contingent},
not \emph{demography-contingent}. This is critical for VRA compliance:
an algorithm that performed worse in states with growing minority
populations would face legal challenge.

\subsubsection{Property 4: Partisan Neutrality}

\begin{quote}
\textit{Without accessing partisan data, the algorithm produces plans
with mean absolute efficiency gap $|EG| = 0.04$---less than half the
$0.08$ mean for enacted plans. The residual imbalance reflects
political geography, not algorithmic choice.}
\end{quote}

Evidence: C.5 documents the $8.3$ percentage-point partisan gap
between algorithmic and enacted plans across 15 competitive states.
The compactness-EG separation (Arizona, Nevada) proves that enacted
partisan bias is not compactness-driven. Five independent fairness
metrics converge on the same substantive conclusion.

Interpretation: The algorithm's partisan neutrality is a structural
property of its information access, not a lucky outcome. Because the
algorithm cannot see partisan data, it cannot gerrymander. The
impossibility defense \citep{b1paper} is validated empirically: the
$|EG|$ of $0.04$ is consistent with the lower bound achievable under
compact districting with Democratic urban concentration, and the
enacted $|EG|$ of $0.08$ exceeds this lower bound, demonstrating
manipulation beyond geographic necessity.

\subsubsection{Property 5: VRA Compliance}

\begin{quote}
\textit{In the 5 states subject to Section 5 preclearance history
(Alabama, Georgia, Louisiana, Mississippi, South Carolina), the
VRA-aligned weighting mode achieves majority-minority district targets
in 4 of 5 states, consistent with the 42\% threshold finding of D.1.}
\end{quote}

Evidence: C.1 finds that majority-minority district counts vary by
$\pm 1$ nationally across a 130$\times$ unit-count range, confirming
that VRA-relevant district structure is resolution-robust. D.1
\citep{d1paper} establishes that the VRA-aligned mode achieves
majority-minority targets in 4 of 5 covered states (the exception is
Louisiana, where district compactness constraints interact with
dispersed minority population geography in ways that require ongoing
study).

Interpretation: VRA compliance is not purely a post-hoc check on
algorithmic output; the algorithm's weight structure can be tuned to
respond to VRA requirements. The 4/5 coverage rate across the most
demanding Section 5 states provides baseline empirical evidence that
the algorithm is suitable for VRA-covered jurisdictions, while the
one failure mode (Louisiana-type dispersed geography) is identified
as a known limitation requiring additional work.

\subsection{The Four Properties as a System}

These four properties interact and reinforce each other in the context
of DIA requirements.

\textbf{Spatial robustness} (P1) establishes that the algorithm's
conclusions are not artifacts of the tract-level implementation choice.
This is a prerequisite for \textbf{temporal stability} (P2): if
changing spatial resolution changed the algorithm's behavior, then
changing spatial building blocks between census years (tract boundaries
are redrawn each decade) could also change behavior, undermining
longitudinal comparisons.

\textbf{Temporal stability} (P2) creates the precondition for
\textbf{demographic neutrality} (P3): because temporal variance is
small and geographically explained, we can rule out the confound that
the algorithm's behavior changes in response to demographic shifts.
If temporal variance were large and unexplained, demographic neutrality
claims would be unsupported.

\textbf{Demographic neutrality} (P3) strengthens the \textbf{partisan
neutrality} (P4) claim. Partisan composition is correlated with
demographic composition; an algorithm that performed systematically
differently for states with high minority populations would face
charges of indirect partisan bias. P3 rules out this mechanism.

Together, the four properties constitute what we call
\textbf{comprehensive algorithmic validity}: the algorithm produces
consistent, stable, demographically neutral, and partisan-neutral
outputs across the full range of conditions under which it would be
deployed.

\subsection{Connection to DIA Reproducibility Requirements}

The DIA's three reproducibility requirements map onto the four
validated properties as follows:

\begin{itemize}
  \item \textbf{R1} (Deterministic from seed): Satisfied by B.7's
    demonstration that CV $<$ 2\% for 48/50 states across 10,000
    seeds, and the DIA seed is statistically indistinguishable from
    optimal.

  \item \textbf{R2} (Invariant to input perturbations): Satisfied by
    P1 (spatial robustness) and P3 (demographic neutrality). The
    algorithm's outputs are stable under the largest practical
    perturbation---a 130$\times$ change in unit count.

  \item \textbf{R3} (Documented behavior across geographies and census
    years): Satisfied by P2 (temporal stability), P3, and P4.
    The C-track provides quantitative characterization of algorithm
    behavior across all 50 states and all three available census years,
    with identified failure modes (high-growth Nevada-style metropolitan
    areas) and their expected magnitudes.
\end{itemize}

\subsection{Limitations of the Validation Program}

The C-track validation has four acknowledged limitations.

First, the block-level analysis (C.1) is complete for 10 of 50 states;
full 50-state block-level validation awaits additional computation.
The 10 states span diverse geographies and the results are consistent,
but the generalization is extrapolated.

Second, C.3 and C.4's temporal stability findings are based on 2 or 3
time points per state. With three census cycles, it is not possible
to distinguish a trend from a level shift. The 2030 census will provide
a critical fourth data point.

Third, the partisan fairness analysis (C.5) is limited to 15
competitive states because single-district states and strongly
non-competitive states do not provide meaningful efficiency gap
estimates. The 15-state sample covers approximately 285 of 435
congressional districts.

Fourth, validation against enacted plans assumes those plans constitute
an appropriate comparison. As reform spreads and enacted plans improve
(C.4's convergence finding), the comparison benchmark itself changes.
Future validation will require a stable reference standard, perhaps
provided by independent redistricting commissions' outputs.
